\chapter{Théorème de Pythagore \& Trigonométrie}
\label{chap:pythagore_trigo}

% ==============================================================================
% 1. OBJECTIFS & COMPÉTENCES DU RÉFÉRENTIEL
% ==============================================================================
\begin{cadreretenu}[Objectifs \& Compétences du DNB Pro]
\begin{itemize}
    \item \textbf{Connaissances} : Triangle rectangle, hypoténuse, théorème direct de Pythagore, réciproque et contraposée, formules de trigonométrie ($\cos, \sin, \tan$).
    \item \textbf{Capacités} :
    \begin{itemize}
        \item \capRealiser\ Calculer la longueur de l'hypoténuse ou d'un côté de l'angle droit avec le théorème de Pythagore.
        \item \capAnalyser\ Démontrer si un triangle est rectangle ou non en utilisant la réciproque ou la contraposée.
        \item \capRealiser\ Calculer la mesure d'un angle aigu à l'aide des touches $\arccos, \arcsin, \arctan$ de la calculatrice.
        \item \capRealiser\ Calculer une longueur inconnue à l'aide des rapports trigonométriques.
        \item \capValider\ Résoudre un problème géométrique de charpente, menuiserie, pente de toit ou accessibilité PMR.
    \end{itemize}
\end{itemize}
\end{cadreretenu}

\vspace{0.3cm}

% ==============================================================================
% 2. COURS ET NOTIONS ESSENTIELLES
% ==============================================================================
\section{Cours : Théorème de Pythagore}

\subsection{Énoncé Direct et Calcul de Longueurs}

\begin{cadreprop}[Théorème de Pythagore]
Dans un triangle $ABC$ \textbf{rectangle en $A$}, le carré de la longueur de l'hypoténuse est égal à la somme des carrés des longueurs des deux autres côtés :
\[ BC^2 = AB^2 + AC^2 \]
\end{cadreprop}

\begin{center}
\begin{tikzpicture}[scale=0.9]
    \coordinate (A) at (0,0);
    \coordinate (B) at (4,0);
    \coordinate (C) at (0,3);

    \draw[fill=colDefBg, draw=colDef, line width=1.2pt] (A) -- (B) node[midway, below] {Côté $= 4\text{ cm}$} -- (C) node[midway, above right] {Hypoténuse $BC = 5\text{ cm}$} -- cycle node[midway, left] {Côté $= 3\text{ cm}$};

    % Angle droit
    \draw[thick, slate800] (0,0.3) -- (0.3,0.3) -- (0.3,0);
    \node[below left] at (A) {$A$};
    \node[below right] at (B) {$B$};
    \node[above left] at (C) {$C$};
\end{tikzpicture}
\end{center}

\begin{cadremethode}[Règles de calcul de longueurs]
\begin{itemize}
    \item \textbf{Pour calculer l'hypoténuse $BC$} :
    $BC^2 = AB^2 + AC^2 \implies BC = \sqrt{AB^2 + AC^2}$.
    \item \textbf{Pour calculer un côté de l'angle droit $AB$} :
    $AB^2 = BC^2 - AC^2 \implies AB = \sqrt{BC^2 - AC^2}$.
\end{itemize}
\end{cadremethode}

\subsection{Réciproque et Contraposée : Démontrer qu'un triangle est rectangle}

\begin{cadreprop}[Réciproque de Pythagore]
Dans un triangle, si le carré du plus grand côté est égal à la somme des carrés des deux autres côtés ($BC^2 = AB^2 + AC^2$), alors ce triangle est \textbf{rectangle} en $A$.
\end{cadreprop}

\begin{cadrepiege}[Rédaction rigoureuse exigée au DNB]
Au Brevet, il est \textbf{strictement interdit} d'écrire l'égalité dès le début ! On doit séparer les calculs :
\begin{enumerate}
    \item On calcule d'une part le carré du plus grand côté : $BC^2 = \dots$
    \item On calcule d'autre part la somme des carrés des deux autres côtés : $AB^2 + AC^2 = \dots$
    \item On compare les deux résultats et on conclut avec la propriété adéquate.
\end{enumerate}
\end{cadrepiege}

% ==============================================================================
% 3. COURS : TRIGONOMÉTRIE DANS LE TRIANGLE RECTANGLE
% ==============================================================================
\section{Cours : Trigonométrie}

\begin{cadredef}[Cosinus, Sinus, Tangente]
Dans un triangle rectangle, pour un angle aigu $\widehat{B}$ :
\begin{itemize}
    \item $\cos(\widehat{B}) = \dfrac{\text{Côté Adjacent}}{\text{Hypoténuse}} = \dfrac{AB}{BC}$
    \item $\sin(\widehat{B}) = \dfrac{\text{Côté Opposé}}{\text{Hypoténuse}} = \dfrac{AC}{BC}$
    \item $\tan(\widehat{B}) = \dfrac{\text{Côté Opposé}}{\text{Côté Adjacent}} = \dfrac{AC}{AB}$
\end{itemize}
\end{cadredef}

\begin{cadreretenu}[Moyen Mnémotechnique Incontournable : \textbf{CAH-SOH-TOA}]
\begin{center}
\textbf{C}osinus = \textbf{A}djacent / \textbf{H}ypoténuse \quad $\vert$ \quad
\textbf{S}inus = \textbf{O}pposé / \textbf{H}ypoténuse \quad $\vert$ \quad
\textbf{T}angente = \textbf{O}pposé / \textbf{A}djacent
\end{center}
\end{cadreretenu}

% ==============================================================================
% 4. AUTOMATISMES & QUESTIONS FLASH
% ==============================================================================
\section{Questions Flash / Automatismes}

\begin{automatisme}[Calculs Rapides \& Règle des Carrés (10 min)]
\noindent
\begin{tabularx}{\linewidth}{X|X}
\textbf{1.} $3^2 + 4^2 = \pointille[1.5cm] \implies \sqrt{\dots} = \pointille[1cm]$ & \textbf{4.} Dans le triangle rectangle, le plus grand côté s'appelle : \pointille[2.5cm] \\
\textbf{2.} $10^2 - 8^2 = \pointille[1.5cm] \implies \sqrt{\dots} = \pointille[1cm]$ & \textbf{5.} $\cos(60\degre) = \pointille[1.5cm]$ (valeur exacte ou décimale) \\
\textbf{3.} Les triplets pythagoriciens usuels : $(3, 4, 5)$ et $(6, 8, \pointille[1cm])$ & \textbf{6.} Sur la calculatrice, pour trouver l'angle dont le cosinus vaut 0,5, on tape : \pointille[2cm] \\
\end{tabularx}
\end{automatisme}

% ==============================================================================
% 5. TRAVAUX DIRIGÉS (TD)
% ==============================================================================
\section{Travaux Dirigés (TD) : Exercices d'Entraînement}

\begin{exercice}[3 pts]{\capRealiser}{Calcul direct de longueur}
Soit $RST$ un triangle rectangle en $S$ tel que $RS = 4{,}8\text{ cm}$ et $ST = 6{,}4\text{ cm}$.
\begin{enumerate}
    \item Faire un schéma à main levée.
    \item Calculer la longueur exacte de l'hypoténuse $RT$.
\end{enumerate}
\lignesreponse[4]
\end{exercice}

\begin{exercice}[3 pts]{\capAnalyser}{Démontrer si un triangle est rectangle}
Un triangle $KLM$ a pour dimensions : $KL = 7\text{ cm}$, $LM = 24\text{ cm}$ et $KM = 25\text{ cm}$.
Le triangle $KLM$ est-il rectangle ? Rédiger avec rigueur.
\lignesreponse[4]
\end{exercice}

\begin{exercice}[4 pts]{\capRealiser}{Calcul d'un angle avec la trigonométrie}
Dans un triangle $ABC$ rectangle en $A$, on donne $AB = 5\text{ cm}$ et $BC = 8\text{ cm}$.
\begin{enumerate}
    \item Écrire la relation trigonométrique donnant $\cos(\widehat{ABC})$.
    \item En déduire la mesure de l'angle $\widehat{ABC}$ arrondie au degré près.
    \item Calculer alors la mesure de l'angle $\widehat{ACB}$.
\end{enumerate}
\lignesreponse[4]
\end{exercice}

% ==============================================================================
% 6. VERS LE BREVET (DNB PRO)
% ==============================================================================
\section{Vers le Brevet (DNB Pro)}

\begin{sujetdnb}[DNB Pro Polynésie][9 pts]{Charpente \& Pente de Toiture}
Un charpentier pose les fermettes d'un comble de maison. La figure ci-dessous représente la coupe d'une ferme de charpente :

\begin{center}
\begin{tikzpicture}[scale=0.9]
    % Triangle principal
    \coordinate (A) at (0,0);
    \coordinate (B) at (8,0);
    \coordinate (S) at (4,2.5);
    \coordinate (H) at (4,0);

    \draw[fill=amber!15, draw=amber!80!black, line width=1.5pt] (A) -- (B) -- (S) -- cycle;
    \draw[dashed, thick, slate700] (S) -- (H) node[midway, right] {Poinçon $SH = 2{,}4\text{ m}$};

    % Angle droit
    \draw[thick] (3.7,0) -- (3.7,0.3) -- (4,0.3);

    \node[below left] at (A) {$A$};
    \node[below right] at (B) {$B$};
    \node[above] at (S) {Sommet $S$};
    \node[below] at (H) {$H$};

    \draw[<->, >=stealth, slate700] (0,-0.4) -- (8,-0.4) node[midway, below] {Entrait $AB = 8\text{ m}$};
    \node at (2,-0.2) [below, font=\scriptsize] {$AH = 4\text{ m}$};
    \node at (6,-0.2) [below, font=\scriptsize] {$HB = 4\text{ m}$};
\end{tikzpicture}
\end{center}

Le triangle $ASB$ est isocèle en $S$. Le point $H$ est le milieu du segment $[AB]$ et le triangle $AHS$ est rectangle en $H$.

\begin{enumerate}
    \item \capSApproprier\ Justifier que la longueur $AH = 4\text{ m}$.
    \item \capRealiser\ Dans le triangle $AHS$ rectangle en $H$, calculer la longueur de l'arbalétrier $[AS]$. Arrondir au centième de mètre.
    \item \capRealiser\ Calculer l'angle de pente du toit $\widehat{SAH}$ formé avec l'horizontale. Arrondir au dixième de degré.
    \item \capValider\ Pour installer des tuiles en terre cuite, le DTU (Document Technique Unifié) exige une pente minimale de $30\degre$. La toiture est-elle conforme à cette réglementation ? Justifier.
\end{enumerate}
\lignesreponse[7]
\end{sujetdnb}

% ==============================================================================
% 7. CORRIGÉS DÉTAILLÉS
% ==============================================================================
\section{Corrigés Détaillés du Chapitre}

\begin{solution}[Corrigé des Exercices et du Sujet DNB]
\textbf{Exercice 1} :
Dans le triangle $RST$ rectangle en $S$, d'après le théorème de Pythagore :
$RT^2 = RS^2 + ST^2 = 4{,}8^2 + 6{,}4^2 = 23{,}04 + 40{,}96 = 64$.
Donc $RT = \sqrt{64} = \repbox{8\text{ cm}}$.

\vspace{2mm}
\textbf{Exercice 2} :
Le plus grand côté est $[KM]$ : $KM^2 = 25^2 = 625$. \\
D'autre part : $KL^2 + LM^2 = 7^2 + 24^2 = 49 + 576 = 625$. \\
On constate que $KM^2 = KL^2 + LM^2$. D'après la réciproque du théorème de Pythagore, le triangle $KLM$ est \textbf{rectangle en $L$}.

\vspace{2mm}
\textbf{Sujet DNB Pro (Charpente)} :
1. $H$ est le milieu de $[AB]$, donc $AH = 8 / 2 = \repbox{4\text{ m}}$. \\
2. Dans le triangle $AHS$ rectangle en $H$, d'après Pythagore : \\
$AS^2 = AH^2 + SH^2 = 4^2 + 2{,}4^2 = 16 + 5{,}76 = 21{,}76 \implies AS = \sqrt{21{,}76} \approx \repbox{4{,}66\text{ m}}$. \\
3. Dans le triangle $AHS$ rectangle en $H$ : $\tan(\widehat{SAH}) = \frac{SH}{AH} = \frac{2{,}4}{4} = 0{,}6$. \\
D'où $\widehat{SAH} = \arctan(0{,}6) \approx \repbox{31{,}0\degre}$. \\
4. L'angle calculé $31{,}0\degre \ge 30\degre$, la pente de la toiture est donc \textbf{pleinement conforme au DTU}.
\end{solution}
